\documentclass[a4paper,11pt]{article}
\usepackage{booktabs}
\usepackage{array}
\usepackage{amsmath}
\usepackage{geometry}
\usepackage{hyperref}
\usepackage{enumitem}
\usepackage{caption}
\geometry{margin=25mm}

\hypersetup{
  colorlinks=true,
  linkcolor=blue,
  urlcolor=blue,
  citecolor=blue
}

\title{A Literature Review on Spiral Unfolding of Polyhedra and Polytopes}
\author{}
\date{April 2026}

\begin{document}
\maketitle
\tableofcontents
\bigskip

% -------------------------------------------------------
\section{Introduction}

The problem of constructing \textbf{nets (unfoldings)} of polyhedra and polytopes is a classical
question dating back to Albrecht D\"{u}rer in the fifteenth century, and at the same time
remains an active research topic in computational geometry and combinatorics.
This review organizes prior work related to Apple-Peel unfolding under the following five themes:
\begin{enumerate}[nosep]
  \item The D\"{u}rer problem and fundamental unfolding methods,
  \item Hamiltonian and Zipper unfoldings,
  \item Spiral unfolding and Apple-Peel unfolding,
  \item Band unfolding,
  \item Unfolding of four-dimensional polytopes.
\end{enumerate}

% -------------------------------------------------------
\section{The D\"{u}rer Problem and Fundamental Unfolding Methods}

\subsection{Origins and Formulation}

The origin of polyhedral nets can be traced to the Renaissance artist \textbf{Albrecht D\"{u}rer}
and his work \textit{Underweysung der Messung} (1525),
which contains hand-drawn nets of solids such as the cube and the regular dodecahedron,
and implicitly raises the question of whether cutting along edges always yields a non-overlapping net.
This question was formulated as a precise mathematical proposition approximately 450 years later
by \textbf{Shephard (1975)}.

\begin{quotation}\noindent
  G. C. Shephard,
  \textit{Convex Polytopes with Convex Nets},
  Mathematical Proceedings of the Cambridge Philosophical Society,
  Vol.~78, pp.~389--403 (1975).
  \url{https://doi.org/10.1017/S0305004100051860}
\end{quotation}

Shephard posed the conjecture that every convex polyhedron has a non-overlapping net
obtainable by cutting along edges (still open today),
and introduced the concept of Hamiltonian unfolding---unfolding along a Hamiltonian edge path.

\subsection{Guaranteed Unfolding Methods}

Two unfolding methods are known that \textbf{guarantee} non-overlap for convex polyhedra.

\begin{quotation}\noindent
  B. Aronov and J. O'Rourke,
  \textit{Nonoverlap of the Star Unfolding},
  Discrete \& Computational Geometry, Vol.~8, No.~3, pp.~219--250 (1992).
  \url{https://doi.org/10.1007/BF02293047}
\end{quotation}

The \textbf{star unfolding} cuts along geodesic shortest paths from a source point to every
vertex, and has been proven to yield a non-overlapping net for any convex polyhedron.
Together with the source unfolding (cutting along the Voronoi diagram), these constitute
the two guaranteed unfolding methods for convex polyhedra.

\subsection{Progress via Affine Transformations}

\begin{quotation}\noindent
  M. Ghomi,
  \textit{Affine Unfoldings of Convex Polyhedra},
  Geometry \& Topology, Vol.~18, No.~5, pp.~3055--3090 (2014).
  arXiv:1305.3231.
\end{quotation}

Ghomi proved that every convex polyhedron in general position is edge-unfoldable after
applying an affine transformation (linear stretching).
This is among the most significant recent advances, showing that there is no combinatorial
obstruction to the D\"{u}rer problem.

\subsection{Comprehensive Reference}

\begin{quotation}\noindent
  E. D. Demaine and J. O'Rourke,
  \textit{Geometric Folding Algorithms: Linkages, Origami, Polyhedra},
  Cambridge University Press (2007).
\end{quotation}

The definitive textbook on computational geometry of folding and unfolding.
It covers all related topics including nets, the D\"{u}rer problem, Cauchy's rigidity theorem,
and Alexandrov's uniqueness theorem, and lists numerous open problems.

% -------------------------------------------------------
\section{Hamiltonian Unfolding and Zipper Unfolding}

The line of work closest in spirit to ``spiral unfolding'' is
\textbf{Hamiltonian unfolding (Zipper unfolding)}.

\begin{quotation}\noindent
  A. Lubiw, E. D. Demaine et al.,
  \textit{Zipper Unfoldings of Polyhedral Complexes},
  Proceedings of CCCG 2010, pp.~219--222.
\end{quotation}

A \textbf{zipper unfolding} is a net obtained by cutting all faces along a single continuous
Hamiltonian edge path (the edge-cutting counterpart of Apple-Peel unfolding).
It has been shown that all Platonic and Archimedean solids admit Hamiltonian (zipper) unfoldings,
making this the direct conceptual predecessor of Apple-Peel unfolding.

\begin{quotation}\noindent
  E. D. Demaine, M. L. Demaine and R. Uehara,
  \textit{Zipper Unfolding of Domes and Prismoids},
  Proceedings of CCCG 2013.
\end{quotation}

For dome-shaped polyhedra, the authors exhibit cases where every Hamiltonian unfolding overlaps
even when the graph is Hamiltonian, while proving that nested prismoids always admit
Hamiltonian unfoldings.

\begin{quotation}\noindent
  J. O'Rourke,
  \textit{Flat Zipper-Unfolding Pairs for Platonic Solids},
  arXiv:1010.2450; CCCG 2010.
\end{quotation}

This paper shows that the tetrahedron, cube, octahedron, and icosahedron each possess a
``zipper pair''---a Hamiltonian edge-path unfolding that refolds into a doubly-covered
parallelogram. No such construction exists for the dodecahedron.

% -------------------------------------------------------
\section{Spiral Unfolding and Apple-Peel Unfolding}

\subsection{Spiral Unfolding (O'Rourke, 2015)}

The most direct predecessor of Apple-Peel unfolding is the following paper by O'Rourke.

\begin{quotation}\noindent
  J. O'Rourke,
  \textit{Spiral Unfoldings of Convex Polyhedra},
  arXiv:1509.00321 (2015).
\end{quotation}

This paper defines a cutting path that covers all faces of a convex polyhedron in a spiral band
(\textbf{spiral unfolding}) and systematically studies whether the resulting net overlaps.
It proves that all 13 Platonic and Archimedean solids admit non-overlapping spiral unfoldings,
taking polyhedra of revolution as the primary class of analysis.
The paper also shows that spiral unfoldings of general convex polyhedra may overlap.
The Apple-Peel unfolding in our project can be seen as a generalization that replaces the
continuous band by a \textbf{discrete greedy selection} of faces.

\subsection{Apple-Peel Unfolding (Yoshino \& Chaidee)}

\begin{quotation}\noindent
  Y. Yoshino and S. Chaidee,
  \textit{Apple Peel Unfolding of Archimedean and Catalan Solids},
  Submitted to Journal of Mathematics and Arts ($\sim$2024).
\end{quotation}

This paper (under review) constructs Apple-Peel unfoldings for all 13 Archimedean solids
and their duals (Catalan solids) and investigates the non-overlap condition.
It is the direct source of the three-dimensional algorithm used in the present project.

\subsection{Prior Work in the Japanese Literature}

\begin{quotation}\noindent
  Unno et al.,
  \textit{On the Spiral Pattern of Apple Peeling and Nets of Regular Polyhedra} [in Japanese],
  \textit{Journal of the Science of Form}, Vol.~27, No.~1, pp.~27--28 (2012).
\end{quotation}

An early Japanese paper discussing the geometric relationship between the spiral pattern
arising when peeling an apple and the nets of regular polyhedra.
It shares the core intuition behind the Apple-Peel concept.

% -------------------------------------------------------
\section{Band Unfolding}

A related line of research is \textbf{band unfolding}, which studies the unfolding of the
portion of a polyhedral surface trapped between two parallel planes (a polyhedral band).

\begin{quotation}\noindent
  G. Aloupis, E. D. Demaine et al.,
  \textit{Edge-Unfolding Nested Polyhedral Bands},
  Computational Geometry: Theory and Applications, Vol.~39, No.~1, pp.~30--42 (2008).
\end{quotation}

A foundational result proving that any nested band---one whose boundary polygon on one side
is contained in the orthogonal projection of the other---is always edge-unfoldable.
This provides a theoretical basis for the ``peel one band at a time'' step in spiral unfolding.

\begin{quotation}\noindent
  M. Radons,
  \textit{Edge-Unfolding Nested Prismatoids},
  Computational Geometry: Theory and Applications (2023).
\end{quotation}

This paper proves that every nested prismatoid (the convex hull of two parallel convex polygons)
is edge-unfoldable, resolving an open problem left by a counterexample in O'Rourke (2007).

% -------------------------------------------------------
\section{Unfolding of Four-Dimensional Polytopes}

\subsection{Combinatorial Enumeration}

\begin{quotation}\noindent
  F. Buekenhout and M. Parker,
  \textit{The Number of Nets of the Regular Convex Polytopes in Dimension $\leq 4$},
  Discrete Mathematics, Vol.~186, pp.~69--94 (1998).
\end{quotation}

A foundational paper classifying and counting all spanning trees of the face-adjacency graph
(up to symmetry equivalence) for every regular convex polytope in dimension $\leq 4$.
Most notably: \textbf{the tesseract (hypercube) has exactly 261 distinct nets}
(a rigorous verification of Turney's 1984 computation).

\subsection{Unfoldability of Regular 4-Polytopes}

\begin{quotation}\noindent
  S. Devadoss and M. Harvey,
  \textit{Unfoldings and Nets of Regular Polytopes},
  Computational Geometry: Theory and Applications (2022).
  arXiv:2111.01359.
\end{quotation}

This paper proves that the $n$-cube, $n$-simplex, and 4-orthoplex are ``all-net''
(every edge unfolding yields a valid, non-overlapping net).
It also demonstrates counterexamples (overlapping nets) for orthoplexes in dimension $\geq 5$
and for the 600-cell, making it an important reference for the unfoldability theory
of the polytopes targeted in our project.

\begin{quotation}\noindent
  S. Devadoss, M. Harvey and Z. Zhang,
  \textit{Visualizing and Unfolding Nets of 4-Polytopes},
  SoCG 2022 (Media Exposition), LIPIcs Vol.~224, Article~67.
\end{quotation}

A media exposition paper presenting an interactive web tool for visualizing nets of
4-dimensional polytopes.
Targeting the 4-simplex, 4-cube, and 4-orthoplex, the tool constructs nets step-by-step
by drawing spanning trees on the dual graph.

\subsection{Hamiltonian Unfolding in Four Dimensions}

\begin{quotation}\noindent
  H. A. Akitaya, R. Samanta et al.,
  \textit{Path-Unfolding the Tesseract},
  Fall Workshop on Computational Geometry (FWCG 2024).
\end{quotation}

This paper enumerates all \textbf{path unfoldings} of the tesseract (nets whose dual graph
is a Hamiltonian path), finding a total of 35,520 such unfoldings.
It is a direct extension of the Zipper/Hamiltonian unfolding framework from three to four
dimensions and shares strong conceptual overlap with the four-dimensional Apple-Peel unfolding
in our project.

\subsection{Hamiltonian Structures on Regular 4-Polytopes}

\begin{quotation}\noindent
  C. H. S\'{e}quin,
  \textit{Symmetrical Hamiltonian Manifolds on Regular 3D and 4D Polytopes},
  Bridges Conference 2005, pp.~463--472.
\end{quotation}

A study of Hamiltonian structures on the edge graphs of regular polytopes,
extended from cycles to manifolds.
Among other results, it shows that the 4-simplex admits a M\"{o}bius-strip-like triangular
covering.
This paper is a useful reference for exploring the connection between spiral unfolding and
Hamiltonian structures in four dimensions.

% -------------------------------------------------------
\section{Discussion: Academic Positioning of Apple-Peel Unfolding}

In light of the prior work surveyed above, the Apple-Peel unfolding can be positioned as follows.

\begin{table}[h]
\centering
\caption{Comparison of Apple-Peel unfolding with prior work}
\label{tab:comparison}
\renewcommand{\arraystretch}{1.3}
\begin{tabular}{lp{5.2cm}p{5.8cm}}
\toprule
Aspect & Prior work & Apple-Peel unfolding \\
\midrule
Cut simplicity
  & Edge cuts only (Zipper); face cuts also possible (general)
  & Order of face selection (cut lines not made explicit) \\
Guarantee of spirality
  & O'Rourke (2015): guaranteed for Platonic/Archimedean
  & Spirality achieved by greedy selection (partially guaranteed) \\
Dependence on symmetry
  & Band unfolding: requires nested structure
  & Higher symmetry $\Rightarrow$ higher success rate (demonstrated on Archimedean solids) \\
Extension to 4D
  & Devadoss (2022): theory; Akitaya (2024): enumeration
  & This project: greedy selection + Darboux frame \\
Algorithm
  & Tabu search (Zawallich 2024), etc.
  & Greedy + one-step backtracking \\
\bottomrule
\end{tabular}
\end{table}

A particularly important distinction is that Apple-Peel unfolding explicitly defines a
\textbf{greedy selection rule} for choosing which face to unfold next.
The formulation of the left-turn condition via the Darboux frame
(using a determinant criterion introduced in this project) is a novel contribution
not found in prior spiral unfolding research.

Regarding the \textbf{extension to four dimensions}: whereas Devadoss \& Harvey (2022) discuss
general unfoldings via spanning trees, our project focuses on the \textbf{selection order}
that produces a spiral, placing it closer in spirit to the path-unfolding framework
of Akitaya et al.\ (2024).
The four-dimensional Apple-Peel unfolding is thus not merely an enumeration of nets,
but a novel attempt to give geometric meaning to ``spirality'' in four dimensions.

% -------------------------------------------------------
\section*{References}
\addcontentsline{toc}{section}{References}

\begin{enumerate}[label={[\arabic*]}, leftmargin=2em]
  \item A. D\"{u}rer, \textit{Underweysung der Messung}, 1525.
  \item G. C. Shephard, Convex Polytopes with Convex Nets,
        \textit{Math.\ Proc.\ Cambridge Philos.\ Soc.}, 78:389--403, 1975.
  \item B. Aronov and J. O'Rourke, Nonoverlap of the Star Unfolding,
        \textit{Discrete \& Comput.\ Geom.}, 8(3):219--250, 1992.
  \item F. Buekenhout and M. Parker, The Number of Nets of the Regular Convex Polytopes
        in Dimension $\leq 4$,
        \textit{Discrete Math.}, 186:69--94, 1998.
  \item T. Biedl, E. D. Demaine et al., Unfolding Some Classes of Orthogonal Polyhedra,
        \textit{Proc.\ CCCG 1998}.
  \item M. Bern, E. D. Demaine et al., Ununfoldable Polyhedra with Convex Faces,
        \textit{Comput.\ Geom.: Theory Appl.}, 2003.
  \item C. H. S\'{e}quin, Symmetrical Hamiltonian Manifolds on Regular 3D and 4D Polytopes,
        \textit{Bridges Conference 2005}, pp.~463--472.
  \item E. D. Demaine and J. O'Rourke, \textit{Geometric Folding Algorithms},
        Cambridge Univ.\ Press, 2007.
  \item G. Aloupis, E. D. Demaine et al., Edge-Unfolding Nested Polyhedral Bands,
        \textit{Comput.\ Geom.: Theory Appl.}, 39(1):30--42, 2008.
  \item A. Lubiw, E. D. Demaine et al., Zipper Unfoldings of Polyhedral Complexes,
        \textit{Proc.\ CCCG 2010}, pp.~219--222.
  \item J. O'Rourke, Flat Zipper-Unfolding Pairs for Platonic Solids,
        arXiv:1010.2450, 2010.
  \item Unno et al., On the Spiral Pattern of Apple Peeling and Nets of Regular Polyhedra
        [in Japanese],
        \textit{J.\ Science of Form}, 27(1):27--28, 2012.
  \item E. D. Demaine, M. L. Demaine and R. Uehara, Zipper Unfolding of Domes and Prismoids,
        \textit{Proc.\ CCCG 2013}.
  \item M. Ghomi, Affine Unfoldings of Convex Polyhedra,
        \textit{Geom.\ \& Topology}, 18(5):3055--3090, 2014.
  \item J. O'Rourke, Spiral Unfoldings of Convex Polyhedra,
        arXiv:1509.00321, 2015.
  \item M. Ghomi, D\"{u}rer's Unfolding Problem for Convex Polyhedra,
        \textit{Notices AMS}, 2018.
  \item S. Devadoss and M. Harvey, Unfoldings and Nets of Regular Polytopes,
        \textit{Comput.\ Geom.: Theory Appl.}, 2022. arXiv:2111.01359.
  \item S. Devadoss, M. Harvey and Z. Zhang, Visualizing and Unfolding Nets of 4-Polytopes,
        \textit{SoCG 2022 (Media Exposition)}, LIPIcs Vol.~224, Article~67.
  \item Y. Yoshino and S. Chaidee, Apple Peel Unfolding of Archimedean and Catalan Solids,
        Submitted to \textit{Journal of Mathematics and Arts}, $\sim$2024.
  \item M. Radons, Edge-Unfolding Nested Prismatoids,
        \textit{Comput.\ Geom.: Theory Appl.}, 2023.
  \item Zawallich, Unfolding Polyhedra via Tabu Search,
        \textit{The Visual Computer}, 2024.
  \item H. A. Akitaya, R. Samanta et al., Path-Unfolding the Tesseract,
        \textit{Fall Workshop on Computational Geometry (FWCG 2024)}.
\end{enumerate}

\bigskip
\noindent\small
Among the references cited in this review, the unpublished/submitted work
(Yoshino \& Chaidee $\sim$2024) is based on limited available information;
details should be confirmed with the authors.
In addition, a Japanese conference presentation by Kaino (2019) on four-dimensional
unfolding nets (Symmetry Society of Japan, Kanazawa) could not be confirmed in
international databases and has therefore been omitted from this review.

\end{document}
